The Conjugate Gradients method, often shortened to CG, is an iterative method for solving linear systems of the form \(Ax=b\), where \(A\) is an \(n\times n\) invertible positive definite symmetric matrix, \(x\) and \(b\) being vectors of size \(n\), without explicit calculating \(A^{-1}\). 
To explain how CG work we need to start with another iterative solver Richardson iterations, which is a uses a simple recursive equation to find an approximation for the value of \(x\)
\begin{equation}\label{eq: Richardson iterations}
    x_{i+1}=x_i+r_i,
\end{equation}

\noindent with \(r_i\) being the residual defined as \(r_i=b-Ax_i\) and a initial \(x_0\). This method work on the condition that \(||I-A||<1\), \cite{Vorst03}, which is for a matrix \(A\in \mathbb{R}^{n\times n}\) the matrix norm defined as

\begin{equation*}
    ||A||=\sup_{y\in \mathbb{R}^n,y\ne 0}\frac{||Ay||}{||y||}.
\end{equation*}

\noindent With a given \(x_0\) we can see that the recursive equation, \eqref{eq: Richardson iterations}, is the sum of all previous residuals and \(x_0\), \cite{Vorst03}
\begin{align*}
    x_{i+1}&=x_{i-1}+r_{i-1}+r_i=x_0+r_0+r_1+\cdots +r_i\\
    &=x_0+\sum_{j=0}^{i}(I-A)^jr_0
\end{align*}

Therefore in the subspace span by \(x_0+\{r_0,Ar_0,\cdots,A^ir_0\}\), where the right part is called the \(i\)th dimensional Krylov subspace generated by \(A\) and \(r_0\), denoted as \(\mathcal{K}^i(A,r_0)\). Which means that Richardson iterations approximation comes from an initial approximation \(x_0\) plus a Krylov subspaces of increasing dimension. We can now almost explain CG as it a Krylov subspace method with Ritz-Galerkin condition for the residuals, which is that \(x_i\) is constructed such that the corresponding residual \(b-Ax_i\) is orthogonal to the current Krylov subspace, \cite{Vorst03}. 
\begin{equation}\label{eq: Ritz-Galerkin}
    b-Ax_i\bot\mathcal{K}^i(A;r_0).
\end{equation}

This can be done by creating an orthonormal basis for the Krylov subspace with \(r_0\) as the initial vector, this is typically done with Arnoldi algorithm, \cite{Saad2003}. This algorithm is a modified version of the standard Gram-Schmidt procedure for creating orthonormal basis for a vector space. The procedure can be seen in Algorithm \ref{alg: Arnoldi}.

\begin{algorithm}[H]
    \caption{Arnoldi algorithm \cite{Vorst03}.}\label{alg: Arnoldi}
    \begin{algorithmic}
        \State \(v_1=\frac{r_0}{||r_0||_2}\)
        \Loop{ \(j=1,\cdots,m-1\)}
            \State \(t=Av_j\)
            \Loop{ \(i=1,\cdots,j\)}
                \State \(h_{i,j}=v_i^Tt\)
                \State \(t=t-h_{i,j}v_i\)
            \EndLoop
            \State \(h_{j+1,j}=||t||_2\)
            \State \(v_{j+1}=\frac{t}{h_{j+1,1}}\)
        \EndLoop
    \end{algorithmic}
\end{algorithm}

With Arnoldi's algorithm we get a relation between the linear systems matrix \(A\), a matrix \(V_m\) created from the basis  vectors \(v_1\) to \(v_m\) as columns and an \(m\times m-1\) upper Hessenberg matrix \(H_{m,m-1}\), which is a matrix with all zero entries below the first subdiagonal, with \(h_{i,j}\) as entries, \(H_{m-1,m-1}\) denotes \(H_{m,m-1}\) without the last row,
\begin{align}
    AV_{m-1}&=V_{m}H_{m,m-1}\label{eq: AV=VH}\\
    &=V_{m-1}H_{m-1,m-1}+h_{m,m-1}v_m e_m^T\\
    V^T_{m-1}AV_{m-1}&=H_{m-1,m-1}\label{eq: VAV=H}
\end{align}

\noindent Note that if \(A\) is symmetric then so is \(H_{m-1,m-1}\) 
\begin{align*}
    H_{m-1,m-1}^T & =(V^T_{m-1}AV_{m-1})^T\\
    &=V^T_{m-1}A^T(V_{m-1}^T)^T\\
    &=V^T_{m-1}AV_{m-1}\\
    &=H_{m-1,m-1}
\end{align*}
and because \(H_{m-1,m-1}\) is a Hessenberg matrix, \(H_{m-1,m-1}\) is a tridiagonal matrix for symmetric \(A\). This is important for the recursive relations of the Conjugate Gradients.

With this orthonormal basis a solution to the linear system would consists of find linear combination, \(y_m\) to satisfy \(x_m=x_0+V_my_m\). With the Ritz-Galerkin condition and \eqref{eq: VAV=H} this can be done with
\begin{align*}
    0&=V_m^T(b-Ax_m)\\
    &=V_m^Tb-V_m^TA(x_0+V_my_m)\\
    &=V_m^T(b-Ax_0)-V_m^TAV_my_m\\
    &=V_m^Tr_0-H_{m,m}y_m
\end{align*}
then by \(r_0=||r_0||_2v_1\) and that \(V_m\) consists of an orthonormal basis \(v_i\), \(V_m^Tr_0=||r_0||_2e_1\), with \(e_1\) being the first canonical unit vector in \(\mathbb{R}^m\), we get
\begin{align*}
    H_{m,m}y_m&=||r_0||_2e_1\\
    y_m&=||r_0||_2H_{m,m}^{-1}e_1.
\end{align*}

With this it can also be shown that the residuals \(r_m\) is proportional to the basis vector \(v_{m+1}\), \cite{Saad2003}. This means that the residuals make up a orthogonal basis for \(\mathcal{K}^m(A;r_0)\) and therefore we can rewrite of the previous equations with the residuals as the basis vectors, \(R_m\), and tridiagonal matrix \(T_{m+1,m}\) defined as \(H_{m+1,m}\)
\begin{align}
    AR_{m}&=R_{m+1}T_{m+1,m}\\
    R_{m}^TAR_{m}&=R_{m}^TR_{m}T_{m,m}\label{eq: RAR=RRT}\\
\intertext{and}
    x_m&=x_0+R_m y_m\label{eq: x=x+Ry}\\
    y_m&=T_{m,m}^{-1}e_1.
\end{align}

With all this we can finally derive the recursion relations for CG. The first thing to see is that \(R_{m}^TR_{m}\) is a diagonal matrix with the squared norms of the residuals on the diagonal, then \eqref{eq: RAR=RRT} shows that \(T_{m,m}\) can be transformed into by rowscaling into symmetric positive definite  (SPD) matrix, because \(A\) is SPD. This means that \(T_{m,m}\) can be LU decomposed without pivoting \cite{Vorst03}.
\begin{equation*}
    T_{m,m}=L_m U_m,
\end{equation*}
with \(L_m\) being lower bidiagonal and \(U_m\) being upper bidiagonal with unit diagonal. Then we can then rewrite \eqref{eq: x=x+Ry}
\begin{equation}\label{eq: x=x+RULe}
    x_m=x_0+(R_mU_m^{-1})(L_m^{-1}e_1).
\end{equation}

By the fact that bidiagonal matrices are nonsingular if all main diagonal elements are nonzero and a closed form exists for the inverse \cite{higham23}, the vector \(\bar{\alpha}:=L_m^{-1}e_1\) can be calculated recursively, with \(\delta_0,\delta_1,\ldots,\delta_m\) on the main diagonal and \(\phi_0,\phi_1,\ldots,\phi_{m-2}\) on the subdiagonal of \(L_m\), as
\begin{align*}
    \alpha_0&=\frac{1}{\delta_0}\\
    \alpha_i&=-\frac{\phi_{i-1}\alpha_{i-1}}{\delta_i}.
\end{align*}

Then for \(P_m := R_m U_m^{-1}\implies R_m =P_m U_m \) with \(P_m\) consists of the vectors \(p_0,\ldots,p_{m-1}\) as columns and \(U_m\) being a upper bidiagonal matrix with unit diagonal and \(\beta_0,\ldots,\beta_{m-2}\) on the superdiagonal, implies that the residuals in the columns of \(R_m\) can be calculated as
\begin{align}
    r_{i}&=\alpha_{i-1}Ap_{i-1}+p_{i}\label{eq: r=bp+p}\\
\intertext{therefore \(p_i\) is defined as}
    p_i&=r_i-\beta_{i-1}p_{i-1}.\label{eq: p=r-bp}
\end{align}

\noindent With these recurve relations \eqref{eq: x=x+RULe} can be written as, \cite{Vorst03},
\begin{equation}\label{eq: CG x_m}
    x_m=x_{m-1}+\alpha_{m-1}p_{m-1}.
\end{equation}

With an initial \(x_0\) and \(p_0:=r_0\) the equations \eqref{eq: CG x_m}, \eqref{eq: r=bp+p}, and \eqref{eq: p=r-bp} defines the recurrence relations for CG. Only missing the values of \(\alpha_i\) and \(\beta_i\)  which can be shown to be equal to, \cite{Vorst03} 
\begin{align}
    \alpha_{i-1}&=\frac{\rho_i}{p_{i-1}^T A p_{i-1}},
\intertext{and}
    \beta_{i-1}&=\frac{\rho_i}{\rho_{i-1}},
\intertext{with}
    \rho_i&=r_i^Tr_i.
\end{align}

With these values of \(\alpha_{i-1}\) and \(\beta_{i-1}\) there is the possibility of zero division. For the latter there is the inner product of \(r_{i-1}\) with it self, \(r_{i-1}^T r_{i-1}\). Because it is an inner product it can only equal zero if \(r_{i-1}\) is the zero vector. This will not happen in implementation, as the standard stopping condition for CG is when the norm of the newest residual is smaller than a given tolerance. As for \(\alpha_{i-1}\) because \(A\) is SPD, \((p_{i-1},p_{i-1})_A=p_{i-1}^T A p_{i-1}\) defines a proper inner product. Therefore as before will not equal zero in implementation.